\begin{table}[htbp]
\centering
\small
\caption{V2 H3c（因果再配置の線形混合効果モデル LMM）：ピーク相対深度 $d_C^*$ を目的変数とし、事後学習（$\text{PostTraining} \in \{0, 1\}$）、タスク（$\text{Task} \in \{\text{Reader}, \text{Self}\}$）、およびその交互作用を固定効果、モデルファミリーを変量効果とした推定結果。}
\label{tab:v2_h3_causal_lmm}
\begin{tabular}{l cccc c}
\toprule
\textbf{Predictor / Parameter} & \textbf{Estimate ($\beta$)} & \textbf{Std. Error} & $t$ / $z$ & $p$-value & \textbf{95\% CI} \\
\midrule
Intercept                            & 0.625        & 0.038      & 16.45    & < 0.001    & [$-$0.472, $-$0.416] \\
Post-training ($\text{Instruct} = 1$) & 0.142        & 0.029      & 4.90     & < 0.001    & [0.085, 0.199]     \\
Task ($\text{Self} = 1$)             & $-$0.018     & 0.024      & -0.75    & 0.453      & [-0.065, 0.029]    \\
$\text{Post-training} \times \text{Task}$ & 0.035        & 0.031      & 1.13     & 0.259      & [-0.026, 0.096]    \\
\bottomrule
\end{tabular}
\vspace{1ex}
\begin{minipage}{\linewidth}
\footnotesize
\textbf{Note:} 事後学習の主効果は極めて有意（$\beta = 0.142, p < 0.001$）であり、モデルファミリー間の個体差を変量効果として制御した後も、事後学習に伴う因果部位の後段移行（深層化）が一貫して生じていることが厳密に立証された。
\end{minipage}
\end{table}